% =============================================================================
% PREAMBLE PROFESIONAL — Electrotecnia Industrial
% Tipografia suave y elegante: formulas, tablas y codigo
% =============================================================================

% --- Fuentes (XeLaTeX): serif profesional para texto, math emparejada -------
\usepackage{fontspec}
\usepackage{unicode-math}
\setmainfont{TeX Gyre Termes}            % texto: serif tipo Times, muy legible
\setsansfont{TeX Gyre Heros}             % sans para titulos/notas
\setmonofont{DejaVu Sans Mono}[Scale=0.85]  % codigo: mono compacto
\setmathfont{TeX Gyre Termes Math}       % formulas: emparejadas con el texto
\setmathfont[range=\mathbb]{Latin Modern Math} % doble trazo en negrita

% --- Microtipografia: suaviza el texto justificado ---------------------------
% Nota: XeTeX solo soporta protrusion (no font expansion)
\usepackage{microtype}
\microtypesetup{protrusion=true, expansion=false}

% --- Paleta de colores suaves ------------------------------------------------
\usepackage{xcolor}   % sin [table]: pandoc ya carga xcolor, evita Option clash
\definecolor{azulmarino}{HTML}{1F4E79}   % titulos y reglas principales
\definecolor{azulmedia}{HTML}{2E74B5}    % acentos
\definecolor{grisfondo}{HTML}{F4F6F8}    % fondo de bloques de codigo
\definecolor{griscabecera}{HTML}{DEE7F2} % cabeceras de tabla
\definecolor{grisborde}{HTML}{B8C4D0}    % bordes suaves
\definecolor{gristexto}{HTML}{333333}    % texto codigo

% --- Espaciado y parrafos ----------------------------------------------------
\usepackage{setspace}
\setstretch{1.12}                        % interlineado suave
\setlength{\parskip}{4pt}
\setlength{\parindent}{1.2em}

% --- Enlaces: discretos, sin cajas ------------------------------------------
\usepackage{hyperref}
\hypersetup{
    colorlinks=true,
    linkcolor=azulmedia,
    citecolor=azulmedia,
    urlcolor=azulmedia,
    pdftitle={Electrotecnia Industrial},
    pdfauthor={KyonX1}
}

% =============================================================================
% TABLAS: elegantes con booktabs + cabecera coloreada
% =============================================================================
\usepackage{booktabs}
\usepackage{tikz}        % portada/portadillas a pagina completa (overlay)
\usepackage{longtable}
\usepackage{array}
\usepackage{colortbl}    % \rowcolor para cabeceras (independiente de xcolor[table])
\usepackage{etoolbox}
\usepackage{needspace}
\usepackage{caption}

% Cabecera con fondo suave y texto en negrita azul
\newcommand{\tablacab}[1]{%
  \arrayrulecolor{azulmarino}%
  \rowcolor{griscabecera}%
  \textbf{\textcolor{azulmarino}{#1}}%
}
% Reglas: gruesa arriba, fina bajo cabecera, gruesa abajo
\renewcommand{\heavyrulewidth}{1.1pt}
\renewcommand{\lightrulewidth}{0.5pt}
\renewcommand{\cmidrulewidth}{0.5pt}
\setlength{\tabcolsep}{7pt}
\setlength{\extrarowheight}{3pt}
\renewcommand{\arraystretch}{1.25}
% (filas alternadas requeririan xcolor[table]; se omiten para evitar clash)

% Captions de tabla: pequenos, negrita el numero, italica el titulo
\captionsetup{font=small, labelfont={bf,color=azulmarino}, textfont=it, skip=6pt}

% =============================================================================
% CODIGO: caja suave con quiebre de linea automatico (arregla desbordes)
% =============================================================================
\usepackage{fvextra}
\usepackage{fancyvrb}
\fvset{
    breaklines=true,              % QUIEBRE de lineas largas (print muy largos)
    breakanywhere=false,
    breaksymbol=\tiny$\hookrightarrow$,
    breakindent=2em,
    fontsize=\footnotesize,
    frame=single,
    framerule=0.3pt,
    rulecolor=grisborde,
    framesep=6pt,
    numbers=left,
    numbersep=8pt,
    numberblanklines=false,
    xleftmargin=12pt,
    xrightmargin=4pt,
    obeytabs=true
}
% Fondo suave para los bloques de codigo de pandoc (entorno Shaded)
\usepackage[most]{tcolorbox}
\BeforeBeginEnvironment{Shaded}{\begin{tcolorbox}[enhanced, sharp corners=all,
    colback=grisfondo, colframe=grisborde, boxrule=0.4pt, left=2pt, right=2pt,
    top=2pt, bottom=2pt, breakable, before skip=8pt, after skip=10pt]}
\AfterEndEnvironment{Shaded}{\end{tcolorbox}}

% =============================================================================
% TITULOS: jerarquia clara con color azul marino
% =============================================================================
\usepackage{titlesec}
% Capitulo (H1): barra azul + titulo grande
\titleformat{\section}
  {\Large\bfseries\color{azulmarino}}{\thesection}{0.8em}{}
\titleformat{\subsection}
  {\large\bfseries\color{azulmedia}}{\thesubsection}{0.7em}{}
\titleformat{\subsubsection}
  {\normalsize\bfseries\color{azulmarino}}{\thesubsubsection}{0.6em}{}
\titlespacing*{\section}{0pt}{16pt plus 4pt}{8pt}
\titlespacing*{\subsection}{0pt}{12pt plus 3pt}{6pt}
\titlespacing*{\subsubsection}{0pt}{10pt plus 2pt}{4pt}

% =============================================================================
% ENCABEZADOS Y PIES DE PAGINA
% =============================================================================
\usepackage{fancyhdr}
\pagestyle{fancy}
\fancyhf{}
\fancyhead[LE]{\small\itshape\textcolor{azulmedia}{\leftmark}}
\fancyhead[RO]{\small\itshape\textcolor{azulmedia}{\rightmark}}
\fancyfoot[LE,RO]{\small\bfseries\color{azulmarino}\thepage}
\fancyfoot[CE,CO]{\small\color{grisborde}Electrotecnia Industrial}
\renewcommand{\headrulewidth}{0.4pt}
\renewcommand{\headrule}{\color{azulmarino}\hrule width\headwidth height 0.4pt}
\renewcommand{\footrulewidth}{0pt}

% --- Tabla de contenidos: numeracion ancha (evita solapamientos) -------------
\usepackage{tocloft}
\setlength{\cftsecnumwidth}{2.8em}
\setlength{\cftsubsecnumwidth}{3.8em}
\setlength{\cftsubsubsecnumwidth}{4.8em}
\renewcommand{\cftsecfont}{\color{azulmarino}}
\renewcommand{\cfttoctitlefont}{\huge\bfseries\color{azulmarino}}

% --- Listas con espaciado agradable ------------------------------------------
\usepackage{enumitem}
\setlist{topsep=4pt, itemsep=2pt, parsep=0pt}

% --- Saltos de pagina controlados --------------------------------------------
\widowpenalty=10000
\clubpenalty=10000
\raggedbottom

% --- Comandos personalizados --------------------------------------------------
\newcommand{\ohm}{\,\Omega}
\newcommand{\uV}{\,\mu V}
\newcommand{\uA}{\,\mu A}
\newcommand{\mA}{\,mA}
\newcommand{\kV}{\,kV}
\newcommand{\kA}{\,kA}

% --- Tamaño de archivo: streams SIN comprimir (PDF legible por cualquier
%     visor, solo ocupa mas espacio en disco; no cambia el aspecto) -----------
% Nota: la compresion de streams queda en el nivel por defecto de xdvipdfmx.
% Las imagenes PNG (portada, portadillas, diagramas, hojas) aportan el peso
% principal del PDF (~11 MB), por lo que no hace falta descomprimir streams.
